\chapter{Data Registration with Distance Maps}\label{chap:distance-maps}

TBC

\section{Background}\label{sec:distance-maps:background}

TBC (ICP  and its variants, etc.)

\subsection{Point cloud registration}\label{sec:distance-maps:background:registration}
One of the most popular approaches for point cloud registration is
the \glslink{icp}{\glsentryfull{icp}} algorithm.
Besl and McKay~\cite{besl-mckay:tpami:1992} address registration of 3D~shapes
represented as point clouds, introducing the standard \gls{icp} algorithm.
The following two steps can summarize this algorithm:
(i) compute correspondences between the source ($z_i$) and target ($p_i$)
cloud points (closest nearest neighbor based on Euclidean distance); and
(ii) estimate the rigid transformation of the source \gls{wrt} target cloud
($\mathbf{X}$) that minimizes a point-to-point distance measure between
corresponding points (see \eqref{eq:distance-maps:background:point-to-point}).
Chen and Medioni~\cite{chen-medioni:icra:1991} exploit the tendency of
most range data to be locally planar and introduce a point-to-plane variant
of \gls{icp}. The proposed method exploits surface normal information ($n_i$).
Instead of minimizing the difference between the projected source point and
the target point, the point-to-plane error metric minimizes
the squared dot product of the surface normal with
the difference between the expected projected source point and
the reference target point
(see \eqref{eq:distance-maps:background:point-to-plane}).
Indeed, Rusinkiewicz and Levoy~\cite{rusinkiewicz-levoy:im:2001} compared
different variants of \gls{icp} with respect to point selection,
data association, weighting of associated pairs, outlier rejection, and
error metric and minimization. The point-to-plane~\cite{chen-medioni:icra:1991}
metric demonstrated significantly better performance than
the point-to-point~\cite{besl-mckay:tpami:1992} formulation.
\begin{equation}\label{eq:distance-maps:background:point-to-point}
\mathbf{X} = \arg \min_\mathbf{X} \sum_i \| \mathbf{X} z_i - p_i \|^2
\end{equation}
\begin{equation}\label{eq:distance-maps:background:point-to-plane}
\mathbf{X} = \arg \min_\mathbf{X} \sum_i \| n_i^T \cdot
\left( \mathbf{X} z_i - p_i \right) \|^2
\end{equation}

Unlike traditional \gls{icp} methods, which rely on explicit point
correspondences, the \glslink{ndt}{\glsentryfull{ndt}}
proposed by Biber and Strasser~\cite{biber-strasser:iros:2003} performs
2D~point cloud registration by representing the environment
with local normal distributions. The 2D~space is discretized into cells.
Each cell is associated with a Gaussian distribution defined by
the mean and covariance of the cell points,
which provide information about the surface those points represent.
Then, Biber and Strasser~\cite{biber-strasser:iros:2003} formulate
2D~point cloud registration as the minimization of an error function
derived from these distributions, optimized using the Gauss--Newton algorithm.
Experimental results showed that \gls{ndt}~\cite{biber-strasser:iros:2003} was
able to map large areas without requiring odometry.
Magnusson~\etal~\cite{magnusson:jfr:2007} later extended \gls{ndt} to
3D~point cloud registration by subdividing space into voxels,
each modelled as a normal distribution estimated from
the points contained within it. The alignment is obtained by
evaluating the likelihood of points from one scan under
the distributions of the other scan.
Compared with point-to-point \gls{icp}~\cite{besl-mckay:tpami:1992} using
kd-tree nearest-neighbor search~\cite{bentley:cacm:1975},
3D-\gls{ndt}~\cite{magnusson:jfr:2007} achieves
lower computation cost by avoiding explicit correspondence search,
while also offering a more compact representation of 3D~point cloud
compared to raw data.

While leveraging
\gls{ndt}~\cite{biber-strasser:iros:2003,magnusson:jfr:2007}-like
representations of the scene, \gls{gicp}~\cite{segal:rss:2009} and
\gls{nicp}~\cite{serafin-grisetti:iros:2015} still retain the \gls{icp}
structure for cloud registration.
\gls{gicp}~\cite{segal:rss:2009} retains the standard \gls{icp}
structure for correspondence estimation, typically using Euclidean
closest-neighbor search accelerated with a kd-tree~\cite{bentley:cacm:1975}.
The main modification lies in the minimization objective function,
introducing a plane-to-plane error metric. \gls{gicp} models
local surface information around each point, estimated from
neighboring points via \gls{pca}. By accounting for the covariance of
the matched points, \gls{gicp} weights residuals based on
local geometric uncertainty. Experiments showed that \gls{gicp} was more robust
to incorrect correspondences and more accurate than
standard point-to-point and point-to-plane \gls{icp}.
\gls{nicp}~\cite{serafin-grisetti:iros:2015} extends \gls{gicp}
not only in the optimization stage but also during correspondence search.
Rather than relying purely on Euclidean nearest-neighbor matching,
\gls{nicp}~\cite{serafin-grisetti:iros:2015} uses a line-of-sight-based
association step and rejects candidate pairs when normals are not well defined,
when point distance exceeds a threshold, when curvature mismatch is too large,
or when the normals are insufficiently aligned. In the optimization step,
\gls{nicp}~\cite{serafin-grisetti:iros:2015} augments each point with
its normal and minimizes a Mahalanobis distance in this 6D space.
In the experiments, \gls{nicp}~\cite{serafin-grisetti:iros:2015} was more
accurate than \gls{gicp}~\cite{segal:rss:2009} and
\gls{ndt}~\cite{magnusson:jfr:2007}
registration algorithms on the evaluated datasets.

Furthermore, Montesano~\etal~\cite{montesano:iros:2005} present
a probabilistic scan matching method for 2D~point cloud registration.
The method~\cite{montesano:iros:2005} builds a set of possible correspondences
using an individual compatibility test based on the Mahalanobis distance
between pairs of points, while accounting for correspondence uncertainty.
Then, Montesano~\etal~\cite{montesano:iros:2005} maximize
the probability of scan alignment over the points correspondence distribution.
The proposed algorithm~\cite{montesano:iros:2005} showed faster convergence
rates than standard \gls{icp}~\cite{besl-mckay:tpami:1992}.
Instead of selecting a single set of correspondences as
Montesano~\etal~\cite{montesano:iros:2005}, Censi~\cite{censi:icra:2006}
considers all possible correspondences and generates alignment hypotheses
through the \gls{ght}, while avoiding explicit point-to-point correspondences.
Censi~\cite{censi:icra:2006} computes alignment within
a fully probabilistic framework that incorporates a motion model and estimates
for registration uncertainty, accounting for surface normals.

Similarly, the GMapping~\cite{grisetti:tro:2007} and
Cartographer~\cite{hess:icra:2016} \gls{slam} algorithms use \gls{mle}
probabilistic formulations to estimate the scan alignment \gls{wrt} the map.
The scan matcher in GMapping~\cite{grisetti:tro:2007} seeks
the most likely robot pose by matching the current observation against
the map built so far. The \gls{mle} formulation estimates
the robot pose given the map, the current scan, and an initial pose guess.
This formulation performs a gradient-descent search over
the observation likelihood function defined on the grid map, based on the
beam-endpoint sensor model~\cite{thrun:book:2005}.
Cartographer~\cite{hess:icra:2016} optimizes the scan pose relative to
the current local submap using a Ceres~\cite{ceres:sw:2023}-based scan matcher.
This matcher estimates the pose that maximizes the occupancy probabilities of
the scan points projected onto the submap. The problem is formulated as
a nonlinear least-squares optimization, using a smooth version of
the occupancy probability values associated with the projected scan points.

Registration based on distance maps:
PerfectMatch~\cite{lauer:robocup:2006,sobreira:jint:2019},
HectorSLAM~\cite{kohlbrecher:ssr:2011}.

\subsection{Distance function representations}\label{sec:distance-maps:background:distance-functions}

truncated signed distance function: 10.1109/ICRA48506.2021.9562069

signed distance function: 10.1109/LRA.2021.3052388

10.1109/ICRA40945.2020.9196695
10.1109/ITSC.2016.7795604
10.1109/IROS.2015.7353633

\subsection{Least-squares optimization}\label{sec:distance-maps:background:least-squares}

\section{Distance Maps}\label{sec:distance-maps}

TBC (unsigned distance map structure,
types of distance map formulations -- only NN, NN + distance,
NN + 1st derivative, NN + 1st and 2nd derivative).

\section{Data Registration}\label{sec:distance-maps:data-registration}

TBC

\subsection{Least-squares on manifolds}

TBC

\subsection{Optimization algorithm}
\subsection{Covariance estimation}

\subsection{Problem formulation -- point-to-point error formulation}

TBC

\subsubsection{State space}
\subsubsection{Measurement space}
\subsubsection{Prediction function}
\subsubsection{Error function}

\subsection{Problem formulation -- point-to-plane error formulation}

TBC

\subsubsection{State space}
\subsubsection{Measurement space}
\subsubsection{Prediction function}
\subsubsection{Error function}

\section{Experimental Results and Discussion}\label{sec:distance-maps:results}

TBC

\subsection{Alignment results}
\subsection{Covariance estimation}
\subsection{Computation performance}

\section{Summary and Conclusions}\label{sec:distance-maps:conclusions}

TBC
